% Overby Industries White Paper
% Source: converted from ledger-bridge.md by tools/md2tex.sh. Edit the Markdown, then regenerate.
% Options for packages loaded elsewhere
\PassOptionsToPackage{unicode}{hyperref}
\PassOptionsToPackage{hyphens}{url}
\PassOptionsToPackage{dvipsnames,svgnames,x11names}{xcolor}
%
\documentclass[
  12pt,
  letterpaper,
]{article}
\usepackage{amsmath,amssymb}
\usepackage{iftex}
\ifPDFTeX
  \usepackage[T1]{fontenc}
  \usepackage[utf8]{inputenc}
  \usepackage{textcomp} % provide euro and other symbols
\else % if luatex or xetex
  \usepackage{unicode-math} % this also loads fontspec
  \defaultfontfeatures{Scale=MatchLowercase}
  \defaultfontfeatures[\rmfamily]{Ligatures=TeX,Scale=1}
\fi
\usepackage{lmodern}
\ifPDFTeX\else
  % xetex/luatex font selection
\fi
% Use upquote if available, for straight quotes in verbatim environments
\IfFileExists{upquote.sty}{\usepackage{upquote}}{}
\IfFileExists{microtype.sty}{% use microtype if available
  \usepackage[]{microtype}
  \UseMicrotypeSet[protrusion]{basicmath} % disable protrusion for tt fonts
}{}
\makeatletter
\@ifundefined{KOMAClassName}{% if non-KOMA class
  \IfFileExists{parskip.sty}{%
    \usepackage{parskip}
  }{% else
    \setlength{\parindent}{0pt}
    \setlength{\parskip}{6pt plus 2pt minus 1pt}}
}{% if KOMA class
  \KOMAoptions{parskip=half}}
\makeatother
\usepackage{xcolor}
\usepackage[margin=1in]{geometry}
\usepackage{longtable,booktabs,array}
\usepackage{calc} % for calculating minipage widths
% Correct order of tables after \paragraph or \subparagraph
\usepackage{etoolbox}
\makeatletter
\patchcmd\longtable{\par}{\if@noskipsec\mbox{}\fi\par}{}{}
\makeatother
% Allow footnotes in longtable head/foot
\IfFileExists{footnotehyper.sty}{\usepackage{footnotehyper}}{\usepackage{footnote}}
\makesavenoteenv{longtable}
\setlength{\emergencystretch}{3em} % prevent overfull lines
\providecommand{\tightlist}{%
  \setlength{\itemsep}{0pt}\setlength{\parskip}{0pt}}
\setcounter{secnumdepth}{-\maxdimen} % remove section numbering
% Overby Industries white paper style, for tools/md2tex.sh.
% Matches the preamble of templates/overby-whitepaper.tex: running header,
% bold section titles, blue links.
\usepackage{fancyhdr}
\usepackage{titlesec}
\usepackage{longtable}
\usepackage{array}

\pagestyle{fancy}
\fancyhf{}
\fancyhead[L]{Overby Industries White Paper}
\fancyhead[R]{\thepage}

\titleformat{\section}{\large\bfseries}{\thesection}{1em}{}
\titleformat{\subsection}{\normalsize\bfseries}{\thesubsection}{1em}{}
\ifLuaTeX
  \usepackage{selnolig}  % disable illegal ligatures
\fi
\IfFileExists{bookmark.sty}{\usepackage{bookmark}}{\usepackage{hyperref}}
\IfFileExists{xurl.sty}{\usepackage{xurl}}{} % add URL line breaks if available
\urlstyle{same}
\hypersetup{
  pdftitle={Overby Industries: A Coordinate Is Not Enough --- Binding Mission and Spending Records to a Tamper-Evident Ledger},
  pdfauthor={Overby Industries R\&D Division},
  colorlinks=true,
  linkcolor={blue},
  filecolor={Maroon},
  citecolor={blue},
  urlcolor={blue},
  pdfcreator={LaTeX via pandoc}}

\title{Overby Industries: A Coordinate Is Not Enough --- Binding Mission
and Spending Records to a Tamper-Evident Ledger}
\usepackage{etoolbox}
\makeatletter
\providecommand{\subtitle}[1]{% add subtitle to \maketitle
  \apptocmd{\@title}{\par {\large #1 \par}}{}{}
}
\makeatother
\subtitle{Version 0.1 --- Draft for Community Review}
\author{Overby Industries R\&D Division}
\date{\today}

\begin{document}
\maketitle

{
\hypersetup{linkcolor=}
\setcounter{tocdepth}{2}
\tableofcontents
}
\hypertarget{abstract}{%
\subsection{Abstract}\label{abstract}}

Two Overby formats, OTF-1 for mission telemetry and OTB-1 for spending,
each carry a field meant to hold proof that a record exists and has not
been altered, and both describe it as a blockchain token. The Overby
Ledger replaces the token with a coordinate into a free, hash-chained,
append-only record: protocol version, sequence number and entry hash.
This paper observes that a coordinate alone is not enough. It proves
that an entry exists and is unaltered. It does not prove the entry is
about the packet that carries it, so a valid token pasted into an
unrelated record still checks out. We describe a bridge that closes the
gap by making each ledger entry commit to a SHA-256 digest of the packet
it records, the circularity that creates and how it is broken, the
canonical form that makes the digest reproducible, and how the ledger's
central protection, that people are committed rather than named, is kept
intact when the packet that points at the ledger is public. We state
what is tested, including a reference for the canonical form written in
independent code, and what is not solved.

\begin{center}\rule{0.5\linewidth}{0.5pt}\end{center}

\hypertarget{the-problem}{%
\subsection{1. The problem}\label{the-problem}}

OTF-1 reports mission events: a piece of debris captured, a load of
regolith refined. OTB-1 reports spending: land bought, a laboratory
built, steel procured. Each has a field, \texttt{ProofOfCleanupToken} in
OTF-1 and \texttt{ProofOfInvestmentToken} in OTB-1, intended to hold
``proof this record exists and has not been altered'', and both were
first described as blockchain tokens.

The Overby Ledger is a free alternative. It is a hash-chained,
append-only record with no token, no consensus and no fee, in which
anyone may hold a copy and verify it independently. It does the token's
job at no cost, and its design documentation argues that the audience
these formats address, regulators and insurers, has little reason to be
reassured by a chain. The two fields stay in the schemas as plain
strings, and what changes is that the value is a ledger coordinate:

\begin{verbatim}
otp:1:8412:9f3c1a4e77b20d51
\end{verbatim}

That is: the coordinate format version, the entry's sequence number, and
the entry's hash as 16 hexadecimal digits. Issuing one costs nothing,
and checking one needs a copy of the ledger and nothing else.

The weakness this paper addresses is simple to state. \textbf{A
coordinate proves an entry exists at that position and has not been
altered since. It says nothing about what the entry is about.} A token
lifted from one record and written into another still checks out,
because the entry it names is perfectly real. For a field whose purpose
is to let a reader trust a record, that is not good enough.

\hypertarget{the-design-in-one-paragraph}{%
\subsection{2. The design in one
paragraph}\label{the-design-in-one-paragraph}}

Each ledger entry carries a SHA-256 digest of the packet it records.
Verification takes the packet in hand, computes its digest, reads the
coordinate from its token field, looks up the entry, checks that the
coordinate is valid against the ledger, and checks that the entry's
digest equals the packet's. A token in a packet that is not the one the
entry recorded fails the second check even though it passes the first.

\hypertarget{the-circle-and-how-it-is-broken}{%
\subsection{3. The circle, and how it is
broken}\label{the-circle-and-how-it-is-broken}}

The obvious design has a problem. The token lives inside the packet. The
token is the hash of the entry. The entry commits to the packet. Taken
literally, the entry would have to contain its own hash.

The way out is that \textbf{the digest is taken over the packet with its
token field removed.} Publishing proceeds in a fixed order: compute the
digest of the token-free packet, append the entry carrying that digest,
learn the new entry's coordinate, and only then write the coordinate
into the packet's token field. Verifying reverses the last step. It
removes the token again before computing the digest, so the digest is
the same before and after stamping. The packet is never hashed with the
value that depends on its own hash.

\hypertarget{the-canonical-form}{%
\subsection{4. The canonical form}\label{the-canonical-form}}

A digest is only useful if two parties get the same one. The same record
may be written with its members in a different order, with whitespace,
or with a timestamp in a different offset, and none of that should
change what was published. If it did, one record could be published
twice, and a check could fail for a reason that has nothing to do with
tampering.

So the digest is taken over a \textbf{canonical form}. The packet is
parsed and validated by the reference decoder for its format. Its token
field is cleared. It is then re-encoded by the reference encoder, which
writes a fixed member order, UTC timestamps, and the shortest number
form that reads back as exactly the same value. Anything the decoder
rejects is not publishable, which makes publication also a validity
gate.

There is an established standard for canonical JSON, RFC 8785 {[}4{]}.
It was not used because it solves a more general problem: it
canonicalises any JSON document. Here the decoder is needed anyway to
validate the packet, and re-encoding the decoded value gives a canonical
form for free, with the schema's own member order. The price is that the
canonical form is defined by the reference implementation, so a verifier
in another language has to implement the written specification rather
than reuse a library.

That is why the specification is written down and independently
reproduced. A short Python script, with no code in common with the C++,
reproduces the form from its written description, and the C++ test suite
is required to match the script's output byte for byte on every example
packet. A change that makes the two disagree fails the build.

\hypertarget{what-is-recorded}{%
\subsection{5. What is recorded}\label{what-is-recorded}}

\begin{longtable}[]{@{}
  >{\raggedright\arraybackslash}p{(\columnwidth - 4\tabcolsep) * \real{0.1406}}
  >{\raggedright\arraybackslash}p{(\columnwidth - 4\tabcolsep) * \real{0.4375}}
  >{\raggedright\arraybackslash}p{(\columnwidth - 4\tabcolsep) * \real{0.4219}}@{}}
\toprule\noalign{}
\begin{minipage}[b]{\linewidth}\raggedright
Ledger field
\end{minipage} & \begin{minipage}[b]{\linewidth}\raggedright
OTF-1 event
\end{minipage} & \begin{minipage}[b]{\linewidth}\raggedright
OTB-1 record
\end{minipage} \\
\midrule\noalign{}
\endhead
\bottomrule\noalign{}
\endlastfoot
Citation & \texttt{OTF-1} & \texttt{OTB-1} \\
Kind & \texttt{otf1.} plus the action, for example
\texttt{otf1.captured} & \texttt{otb1.} plus the type, for example
\texttt{otb1.land} \\
Subject & The pod, named. A machine is not a person & The vendor, named
or committed (section 6) \\
Standing & Determined & Determined \\
Time & The packet's timestamp in milliseconds since the Unix epoch & The
same \\
Detail & \texttt{sha256:} and 64 hex digits, then a human-readable
summary & The same \\
\end{longtable}

The summary exists so a person reading a diff of a mirrored ledger can
see what an entry is about, for example ``captured by pod A1, mission
orbital-demo-1, task T-0042''. It is never trusted: verification reads
only the digest.

Only OTF-1 event packets are published. A status, heartbeat, assignment
or acknowledgement is bookkeeping between a vehicle and its coordinator,
not the public record, and is refused.

\hypertarget{keeping-the-asymmetry}{%
\subsection{6. Keeping the asymmetry}\label{keeping-the-asymmetry}}

The ledger's central design decision is that disclosure is asymmetric.
Institutions are named and fully linkable, because pattern over an
institution is what a transparency record exists to reveal. People are
represented by a keyed commitment that cannot be linked across their own
entries or to their identity without a key held separately, because
pattern over a person is what enables targeting. A register of everyone
ever accused of anything, searchable by name, is how a transparency
system becomes an instrument of the harm it was built to surface.

An OTB-1 record names a vendor. The ledger's design distinguishes a
legal entity, which is named, from a natural person such as a sole
trader, which is committed. The transaction itself stays public. What is
hidden is only the counterparty's identity, and only when the
counterparty is a person.

A bridge could defeat this without any bug in the ledger. If the ledger
committed a person while the packet that points at it named them, nobody
would have been protected: the name would sit in public next to the
coordinate. So the bridge enforces the rule on both sides.

\begin{enumerate}
\def\labelenumi{\arabic{enumi}.}
\tightlist
\item
  \textbf{The caller must classify the vendor.} An OTB-1 record cannot
  say whether its vendor is a legal entity. The bridge takes the
  classification as an argument and treats an unclassified vendor as a
  person. That falls protective for the ledger's own reason: wrongly
  naming a person cannot be undone, and wrongly protecting a company
  only means somebody has to ask.
\item
  \textbf{For anyone not established to be a legal entity,} the ledger
  commits the identity and keeps no trace of it, and the packet the
  bridge returns has the vendor replaced by a placeholder \emph{before}
  the digest is taken. The digest, the stamped record and the ledger
  therefore agree on what was published.
\end{enumerate}

The bridge returns the withheld form and nothing else. The original,
un-withheld record is not verifiable: its digest is not the one the
ledger holds. The tests check that no trace of a withheld name reaches
the ledger's canonical text or its JSON.

What this does not do is also stated. The bridge inspects only the
vendor field. A person's name typed into a free-text field such as the
item description is not found. The classification is a decision made by
the caller and an explicit one, so classifying the same vendor
differently is a different record. And the rule is the design's, not a
legal conclusion: financial records carry statutory retention that in
many jurisdictions overrides erasure, and this needs review by someone
qualified before an operations ledger publishes.

\hypertarget{verification}{%
\subsection{7. Verification}\label{verification}}

Verification reports what is wrong rather than only that something is:

\begin{longtable}[]{@{}
  >{\raggedright\arraybackslash}p{(\columnwidth - 2\tabcolsep) * \real{0.2558}}
  >{\raggedright\arraybackslash}p{(\columnwidth - 2\tabcolsep) * \real{0.7442}}@{}}
\toprule\noalign{}
\begin{minipage}[b]{\linewidth}\raggedright
Verdict
\end{minipage} & \begin{minipage}[b]{\linewidth}\raggedright
Meaning
\end{minipage} \\
\midrule\noalign{}
\endhead
\bottomrule\noalign{}
\endlastfoot
Verified & The coordinate checks out and its entry commits to this
packet \\
Not a packet & The text is not valid OTF-1 or OTB-1 \\
Not publishable & Valid OTF-1, but not an event, so nothing would have
been recorded \\
No token & The token field is absent \\
Foreign token & The field holds something that is not a coordinate, such
as a legacy blockchain token. It cannot be checked here \\
Malformed token & It looks like a coordinate and is not one \\
Token invalid & The coordinate fails against this ledger: no such entry
(a copy that is merely behind), a different hash (altered, or issued by
another ledger), or a broken chain \\
Entry mismatch & The coordinate is valid but its entry is about
something else: the packet was altered, or the token was lifted \\
\end{longtable}

Two properties are worth stating. A ledger that is merely behind a
token, a copy taken before the entry was published, reports a missing
entry rather than a forgery, because they are different problems. And a
break in the chain after an entry does not invalidate an earlier token:
an entry is vouched for by what precedes it.

\hypertarget{existing-tokens-and-publishing-twice}{%
\subsection{8. Existing tokens, and publishing
twice}\label{existing-tokens-and-publishing-twice}}

A packet may already carry something in its token field. The bridge will
not overwrite it unless told to, because a token that was meaningful
would be lost silently. The example packets in both format repositories
all carry legacy tokens, so their tests set the option. A packet
carrying a different coordinate from the one the ledger holds for it is
refused for the same reason.

Publishing the same record twice appends it once, however it is spelled
and whether or not it has already been stamped, because the digest
ignores the token. A bridge over a ledger loaded from a copy rebuilds
its index from the entries, so a record published in an earlier session
is found. A ledger that does not say which ledger it is cannot be
published into.

\hypertarget{evaluation}{%
\subsection{9. Evaluation}\label{evaluation}}

The tests are adversarial by design, like the ledger's own. Anyone can
show that a record published and then checked comes back verified. The
tests try every way of presenting one that should not: an altered
packet, a token lifted from another record, a coordinate for an entry
that does not exist, a coordinate issued by a different ledger, a ledger
rewritten after publication, an entry removed from the middle, and a
vendor that must not be named.

SHA-256 is checked against the standard's test vectors {[}1{]} and
against an independent implementation at every padding boundary. The
canonical form and digest are required to match the independent
reference byte for byte. To test the tests, the code was deliberately
broken in fourteen ways, among them a digest that includes the token, a
verifier that skips the comparison, a person named in the packet or the
ledger, a token silently overwritten and an off-by-one in the hash's
padding, and each was caught by a clean assertion failure. That shows
the suite would notice such a regression. It does not show the design is
complete.

\hypertarget{limits}{%
\subsection{10. Limits}\label{limits}}

\begin{itemize}
\tightlist
\item
  \textbf{The chain hash is still FNV-1a, 64 bits.} The digest binding a
  packet to its entry is SHA-256 and does not inherit that weakness. The
  entry itself, and so the coordinate, does. That is fine against
  accident and casual tampering and not sufficient against a determined
  adversary, as the ledger's own documentation says, and a publisher
  with real adversaries should substitute a stronger hash.
\item
  \textbf{A coordinate names no ledger.} Checking it against the wrong
  ledger fails with a hash mismatch, because the instance is inside
  every entry's hash, but a reader is not told which ledger to look in.
\item
  \textbf{Commitments are only as private as their key,} and the ledger
  has no way to notice that the key was never set.
\item
  \textbf{A person is found by classification, not detection.}
\item
  \textbf{It records what is handed to it.} Nothing here stops a
  publisher choosing not to publish an event. That is the job of the
  ledger's gap-free sequence numbers and of distributed copies.
\item
  \textbf{Verification is linear.} Checking a coordinate walks the
  chain.
\end{itemize}

\hypertarget{related-work}{%
\subsection{11. Related work}\label{related-work}}

Tamper-evident logging has a body of work behind it {[}6{]}, and
Certificate Transparency {[}5{]} made one design, a Merkle tree with
compact inclusion proofs, widely deployed. The ledger described here is
a hash chain, which is simpler and gives linear verification, where a
Merkle tree would give logarithmic inclusion proofs a verifier could
check without holding the whole ledger. Moving to one is an open
question below. Canonical JSON is the subject of RFC 8785 {[}4{]},
discussed in section 4. The packets use JSON {[}2{]} and RFC 3339
timestamps {[}3{]}.

\hypertarget{open-questions}{%
\subsection{12. Open questions}\label{open-questions}}

\begin{itemize}
\tightlist
\item
  \textbf{A Merkle structure.} Whether compact inclusion proofs are
  worth the added complexity, so that a reader can check one coordinate
  without holding the whole ledger.
\item
  \textbf{A stronger entry hash,} and whether the coordinate should
  carry a longer one.
\item
  \textbf{The instance in the coordinate,} so a reader is told which
  ledger to consult.
\item
  \textbf{A signature layer} binding a ledger to a publisher's key, and
  where that key is held.
\item
  \textbf{The classification rule,} and who may review it before an
  operations ledger publishes.
\end{itemize}

\hypertarget{references}{%
\subsection{References}\label{references}}

\begin{enumerate}
\def\labelenumi{\arabic{enumi}.}
\tightlist
\item
  National Institute of Standards and Technology, ``Secure Hash Standard
  (SHS),'' FIPS PUB 180-4, 2015.
\item
  T. Bray (ed.), ``The JavaScript Object Notation (JSON) Data
  Interchange Format,'' RFC 8259, 2017.
\item
  G. Klyne and C. Newman, ``Date and Time on the Internet: Timestamps,''
  RFC 3339, 2002.
\item
  A. Rundgren, B. Jordan and S. Erdtman, ``JSON Canonicalization Scheme
  (JCS),'' RFC 8785, 2020.
\item
  B. Laurie, A. Langley and E. Kasper, ``Certificate Transparency,'' RFC
  6962, 2013.
\item
  S. A. Crosby and D. S. Wallach, ``Efficient data structures for
  tamper-evident logging,'' \emph{Proceedings of the 18th USENIX
  Security Symposium}, 2009.
\item
  Overby Industries, \emph{overby-transparency-protocol}, the Overby
  Ledger and its bridge:
  \url{https://github.com/Overby-Industries/overby-transparency-protocol}
  (see \texttt{docs/LEDGER.md} and \texttt{docs/BRIDGE.md}).
\item
  Overby Industries, \emph{telemetry-protocol} (OTF-1):
  \url{https://github.com/Overby-Industries/telemetry-protocol} and
  \emph{transparency-protocol} (OTB-1):
  \url{https://github.com/Overby-Industries/transparency-protocol}.
\end{enumerate}

\end{document}
